\section{问题2：整个烘干过程}

\subsection{模型建立}

问题2 求解\textbf{整个烘干过程}，从 $t=0$ 起统一采用附录3 的变物性经验公式：
\begin{equation}
  \rho(C)c_p(C)\frac{\partial T}{\partial t}
  =\frac{1}{r}\frac{\partial}{\partial r}\!\left(k(C)\,r\frac{\partial T}{\partial r}\right),
  \quad
  \frac{\partial C}{\partial t}
  =\frac{1}{r}\frac{\partial}{\partial r}\!\left(D(C,T)\,r\frac{\partial C}{\partial r}\right),
\end{equation}
\begin{equation}
  D=2.4\times10^{-3}\exp\!\left(-\frac{0.45}{C}\right)
    \exp\!\left[-\frac{3850}{T+273.15}\right]\quad[\mathrm{m^2/s}] .
\end{equation}

\textbf{题面明确``为简化问题，相关经验公式统一采用附录3''}，
故从 $t=0$ 起就用附录3，\textbf{不存在阶段切换}，也不需要``切换处保持状态连续''的处理。

\textbf{双向耦合}：$D$ 依赖 $T$，水分场不再独立，$T$ 与 $C$ 必须\textbf{联立求解}。

\begin{figure}[H]
  \centering
  \includegraphics[width=0.75\textwidth]{FIG-16_dimensionless.png}
  \caption{无量纲数分析。(a) $Bi_m$ 由 1.19 升到 47.7，\textbf{过程中跨越 1}；
           (b) 传质特征时间比传热大 1$\sim$2 个数量级}
  \label{fig:FIG16}
\end{figure}

\begin{figure}[H]
  \centering
  \includegraphics[width=0.78\textwidth]{FIG-18_D_contour.png}
  \caption{附录2/3/4 的扩散系数 $D(C,T)$ 等值线与本文工况轨迹（三者色标范围相同）}
  \label{fig:FIG18}
\end{figure}

\subsection{边界外推与敏感性}

长时环境边界的处理方案见 \S5.3。本节做情景扫描检验其影响：

\begin{table}[H]
  \centering
  \caption{环境外推情景对 $\tstar$ 的影响}
  \label{tab:envsens}
  \footnotesize
  \begin{tabular}{@{}lccc@{}}
    \toprule
    情景 & $\tstar$ / h & $\Delta$ vs 基线 & 相对 \\
    \midrule
    基线（$t_{pre}=14400$ s、$T_\infty=50.0017\ ^\circ$C、$\Cinf=0.049998$） & 57.49 & --- & --- \\
    $t_{pre}=9600$ s & 57.49 & $\approx0$ & $\approx0$ \\
    $T_\infty=49.8\ ^\circ$C & 57.87 & $+0.39$ h & $+0.67\%$ \\
    $T_\infty=50.3\ ^\circ$C & 56.95 & $-0.54$ h & $-0.94\%$ \\
    $\Cinf=0.0495$ & 57.48 & $-0.01$ h & $-0.02\%$ \\
    $\Cinf=0.0505$ & 57.50 & $+0.01$ h & $+0.02\%$ \\
    \bottomrule
  \end{tabular}
\end{table}

\textbf{由环境外推得到的情景区间}：$\tstar\in[56.95,57.87]$ h $=[2.373,2.411]$ 天
（$\pm0.8\%$）。$t_{pre}$ 的取值对结果\textbf{没有影响}——
因为 PCHIP 数据段与平台值本身是连续过渡的。

\textbf{环境平台窗口的选择影响}：本文取 $t\ge9600$ s 段（$n=81$）均值。
若改取 $t\ge12400$ s 段，$\tstar$ 会从 $57.4889$ h 变到 $57.4747$ h，
跨度约 $0.03$ h。这属于\textbf{建模约定不同}，不是模型或算法误差。
完整的口径对照见 \S\ref{sec:tstar-precision}。

\subsection{数值方法}

问题2 在问题1 内核基础上做两项扩展：

\textbf{（1）变物性走守恒形式}。面导热系数 $k$ 用调和平均、
热容 $(\rho c_p)$ 用单元值，\textbf{二者不可合并}。
若令 $\Gamma=k/(\rho c_p)$ 再走原路径，等于把两者在同一个面上一起做调和平均，
$\rho c_p$ 沿 $r$ 变化数倍时会产生明显偏差。
守恒形式的能量方程应写成
\begin{equation}
  \left(\rho c_p\right)_i V_i\frac{T_i^{n+1}-T_i^{n}}{\Delta t}
  =\sum_{j} G_{ij}\left(T_j^{n+1}-T_i^{n+1}\right)
  +\ \text{边界项},
\end{equation}
即 $\rho c_p$ \textbf{整体}除以 $(\rho c_p)_i V_i$，而不是把 $k$ 与 $c_p$ 分别处理。

\textbf{（2）块 Picard 用 Gauss--Seidel 序}。$\rho c_p$、$k$ 对 $C$ 的依赖是单向的，
用最新 $C$ 可显著加快收敛且不破坏任何子问题的 M-矩阵性质。

\textbf{残差判据不可省}：只查更新量会落入``更新量变小但残差停滞''的陷阱；
本文的\textbf{双重判据}对 $T$、$C$ 两个场同时成立才算收敛。

\textbf{网格与时间}：生产网格 $N=200$、$\gamma=1.5$，$\theta=1$（Backward Euler），
自适应步长，输出时刻精确落点。3 h 正式版接受步 12652 次、拒绝步 4 次，
平均 3.00 次 Picard 迭代。

\subsection{结果}

\begin{figure}[H]
  \centering
  \includegraphics[width=0.78\textwidth]{G5_p2_radial.png}
  \caption{问题2 前 3 h 的径向分布}
  \label{fig:G5}
\end{figure}

\textbf{温度场}：前 3 h 内整体升温到接近烘房温度；
$t=3$ h 时中心 $49.8494\ ^\circ\mathrm{C}$、表面 $49.9666\ ^\circ\mathrm{C}$，
表面略高于中心，符合热风加热过程。

\textbf{含水率场}：前 3 h 内含水率显著下降但远未达终点；
$t=3$ h 时中心 $1.7662$ kg/kg、表面 $1.0081$ kg/kg，表面先干、中心后干。

\begin{table}[H]
  \centering
  \caption{3 小时内药材的温度（$^\circ$C）}
  \label{tab:p2T}
  \footnotesize
  % 数据源：甲day2/results/M0/result2.xlsx
  \begin{tabular}{@{}lrrrrr@{}}
    \toprule
    时间/h & 0 cm & 0.5 cm & 1 cm & 1.5 cm & 2 cm \\
    \midrule
    0.5 & 32.1902 & 32.3830 & 32.9666 & 33.9613 & 35.4136 \\
    1 & 40.3821 & 40.5541 & 41.0608 & 41.8783 & 43.0000 \\
    1.5 & 45.8467 & 45.9348 & 46.1929 & 46.6052 & 47.1395 \\
    2 & 48.4500 & 48.4879 & 48.5982 & 48.7734 & 49.0032 \\
    2.5 & 49.4668 & 49.4791 & 49.5135 & 49.5645 & 49.6617 \\
    3 & 49.8494 & 49.8552 & 49.8746 & 49.9100 & 49.9666 \\
    \bottomrule
  \end{tabular}
\end{table}

\begin{table}[H]
  \centering
  \caption{3 小时内药材的水分浓度（kg/kg）}
  \label{tab:p2C}
  \footnotesize
  % 数据源：同上
  \begin{tabular}{@{}lrrrrr@{}}
    \toprule
    时间/h & 0 cm & 0.5 cm & 1 cm & 1.5 cm & 2 cm \\
    \midrule
    0.5 & 2.5499 & 2.5489 & 2.5255 & 2.3257 & 1.6486 \\
    1 & 2.5257 & 2.4948 & 2.3578 & 2.0230 & 1.4711 \\
    1.5 & 2.3860 & 2.3256 & 2.1344 & 1.8020 & 1.3476 \\
    2 & 2.1708 & 2.1084 & 1.9236 & 1.6259 & 1.2311 \\
    2.5 & 1.9567 & 1.9006 & 1.7360 & 1.4720 & 1.1166 \\
    3 & 1.7662 & 1.7165 & 1.5702 & 1.3333 & 1.0081 \\
    \bottomrule
  \end{tabular}
\end{table}


完整结果见 \texttt{result2.xlsx}。

\subsection{与问题1 的差异分析}

问题2 相比问题1 有三项扩展：\textbf{变物性}（$D$ 含 Arrhenius 温度依赖，
方程强非线性）、\textbf{双向耦合}（$T$ 与 $C$ 必须联立）、
\textbf{全流程边界}（跨越 2--3 天）。

在 $t=1800$ s 处，双向耦合使中心温度降低约 $\mathbf{1.386\ ^\circ C}$
（问题1 为 $33.5763\ ^\circ\mathrm{C}$；问题2 为 $32.1902\ ^\circ\mathrm{C}$）。
\textbf{这属于模型差异分析，不是两套实现必须相等的正确性验证。}

\textbf{差异方向值得注意}：表面上附录3 的 $D$ 在 $30\ ^\circ\mathrm{C}$ 附近
反而略大于附录2（$D_{\text{app3}}(2.55,30\ ^\circ\mathrm{C})=5.64\times10^{-9}$
对 $D_{\text{app2}}(2.55)=4.94\times10^{-9}$），\textbf{本应干燥更快}；
但实测问题2 \textbf{更冷、更湿}。原因正是\textbf{双向耦合}：

\begin{itemize}[leftmargin=1.6em, itemsep=2pt]
  \item 附录3 的 $\rho c_p$ 更大、$k$ 更高，热扩散率\textbf{低 14\%} $\to$ 升温更慢；
  \item 而 $D$ 含 Arrhenius 因子 $\exp(-3850/T_K)$（约 $4\%/\mathrm{K}$），
        温度偏低又\textbf{反馈}压低扩散系数。
\end{itemize}

这正是``问题1 可以先水后温、问题2 必须联立''的直接数值体现。

作为对照，M1 在 M0 基础上开启蒸发潜热后，前 3 h 中心温度额外降低
$18.06$ K；M1 为\textbf{条件性扩展，不改变问题2 的正式答案}（见 \S\ref{sec:M1}）。

\subsection{特征时间与物性公式的定量比较}

取初始状态（$T=28\ ^\circ\mathrm{C}$，$C=2.55$ kg/kg）估算特征时间：
\begin{equation}
  \tau_T=\frac{R_0^{2}}{\alpha}\approx2761.9\ \mathrm{s}\approx46.0\ \mathrm{min},
  \qquad
  \tau_m=\frac{R_0^{2}}{D}\approx7.09\times10^{4}\ \mathrm{s}\approx19.7\ \mathrm{h}.
\end{equation}
因此前 3 h 温度已接近平衡，水分仍在缓慢扩散，整个烘干过程由水分扩散主导。

\textbf{关于``265 倍''的说明}：本文按附录公式计算得到附录2、附录3、附录4 的
$D(C)$ 在水分浓度方向上的变化幅度分别为 \textbf{266.2 倍}、\textbf{16.8 倍}、
\textbf{6.6 倍}。\textbf{题面并未出现``265 倍''这一比值}，
该数是本文根据附录数据计算的结果，引用时须写成``本文计算得到''，
\textbf{不得}写成``题目给出''。

\subsection{完整过程的物理本质：恒温干燥}

由时间尺度分析可知，传质特征时间随工况变化很大：

\begin{table}[H]
  \centering
  \footnotesize
  \begin{tabular}{@{}lcc@{}}
    \toprule
    工况 & $\tau_m=R_0^2/D$ / h & $\tau_m/\tau_T$ \\
    \midrule
    $T=28\ ^\circ$C，$C=2.55$ & 19.7 & $\approx26$ \\
    $T=50\ ^\circ$C，$C=2.55$ & 8.2 & $\approx11$ \\
    $T=50\ ^\circ$C，$C=0.15$ & 139 & $\approx181$ \\
    \bottomrule
  \end{tabular}
\end{table}

约 4$\sim$5 个传热时间常数（$\approx3.5$ h）内温度场就基本平衡，
之后是长达约 \textbf{54 h 的恒温干燥}，温度不再是限制因素。
\textbf{整个``2--3 天''的时长由水分扩散单独决定。}
这解释了``预热平衡 $+$ 恒温干燥''两阶段的自然划分，
也说明恒温段的边界设定对 $\tstar$ 的影响远大于对温度场的影响。
